%% toc: P3 40 40.5
\section{Proofs for closure constructions and the dynamic view}
\label{sec:annex-p03-w03-p03-068-closure}

Throughout, the conventions of \cref{def:p03-w03-p03-068-standing} are in
force: \(\alpha\in(0,1)\) is irrational unless a statement says otherwise,
\(t\geq0\) is an integer, and \(e_G(N;S)\) counts the edges of \(G\) with at
least one end in \(N\) and both ends in \(S\cup N\).  Rooted graphs and their
parameters are as in
\cref{def:p03-w03-p03-065-rooted-graph,def:p03-w03-p03-065-dense-sparse}.
An edge of a rooted graph \((S,H)\) is \emph{counted} if it does not join two
roots; the parameter \(e\) of \((S,H)\) is the number of counted edges.

\begin{annexitem}{annex:p03-w03-p03-068-rigid-local}{thm:p03-w03-p03-068-rigid-local}{Rigidity is a property of the induced graph}
\begin{proof}
\emph{(2) and (3) are the same condition.}  The nailextensions of the
rooted graph \((X,G[X\cup Y])\) are the rooted graphs \((S,G[X\cup Y])\) with
\(X\subseteq S\subsetneq X\cup Y\).  The nonroots of \((S,G[X\cup Y])\) are
the vertices of \(Y\setminus S\), and its counted edges are the edges of
\(G[X\cup Y]\) with at least one end outside \(S\), that is, with at least
one end in \(Y\setminus S\) and both ends in \(S\cup(Y\setminus S)=X\cup Y\).
Their number is \(e_G(Y\setminus S;S)\).  So \((S,G[X\cup Y])\) is dense
exactly when \(|Y\setminus S|-e_G(Y\setminus S;S)\alpha<0\), and rigidity of
\((X,G[X\cup Y])\) is the condition defining ``\(Y\) is rigid over \(X\)''.

\emph{(2) implies (1).}  Give \(G[X\cup Y]\) the root labels
\(x_1,\ldots,x_r\) and nonroot labels \(b_1,\ldots,b_v\).  Then
\(b_1,\ldots,b_v\) form an \((X,G[X\cup Y])\) extension over
\(x_1,\ldots,x_r\), and this rooted graph is rigid.

\emph{(1) implies (3).}  Let \(b_1,\ldots,b_v\) form an \((R,H)\) extension
over \(x_1,\ldots,x_r\) with \((R,H)\) rigid, and let \(\iota\colon V(H)\to
X\cup Y\) send the root \(x_i\) of \(H\) to \(x_i\) and the nonroot \(y_j\) to
\(b_j\); it is a bijection.  Let \(X\subseteq S\subsetneq X\cup Y\) and
\(S_H=\iota^{-1}(S)\), so \(R\subseteq S_H\subsetneq V(H)\).  The nailextension
\((S_H,H)\) has the \(|Y\setminus S|\) nonroots \(\iota^{-1}(Y\setminus S)\).
Each of its counted edges has an end outside \(S_H\), which is a nonroot of
\(H\), so it is an edge of the form \(y_iy_j\) or \(x_iy_j\); by the definition
of an extension its image under \(\iota\) is an edge of \(G\), with an end in
\(Y\setminus S\) and both ends in \(X\cup Y\).  Distinct edges have distinct
images, so \(e_G(Y\setminus S;S)\) is at least the parameter \(e\) of
\((S_H,H)\).  Since \((S_H,H)\) is dense,
\[
  |Y\setminus S|-e_G(Y\setminus S;S)\alpha\leq|Y\setminus S|-e\alpha<0 .
\]
As \(S\) was arbitrary, \(Y\) is rigid over \(X\).

\emph{Locality.}  For \(X\subseteq S\subsetneq X\cup Y\), every edge counted
by \(e_G(Y\setminus S;S)\) has both ends in \(X\cup Y\subseteq W\), so it is an
edge of \(G[W]\), and conversely every edge of \(G[W]\) is an edge of \(G\).
Hence \(e_G(Y\setminus S;S)=e_{G[W]}(Y\setminus S;S)\) for every such \(S\), and
the defining inequalities of rigidity coincide in \(G\) and \(G[W]\).  This
argument does not use finiteness of \(X\).
\end{proof}
\end{annexitem}

\begin{annexitem}{annex:p03-w03-p03-068-move-base}{thm:p03-w03-p03-068-move-base}{Moving to a larger base}
\begin{proof}
\(Y'\) is finite and nonempty, disjoint from \(B\), and \(|Y'|\leq|Y|\).  Let
\(S\) satisfy \(B\subseteq S\subsetneq B\cup Y'\) and put
\(N=(B\cup Y')\setminus S\).  Since \(B\subseteq S\), \(N=Y'\setminus S\), and
\(N\neq\emptyset\).  Put \(S_0=(X\cup Y)\setminus N\).  Since \(N\subseteq
Y'\subseteq Y\) and \(Y\cap X=\emptyset\), we have \(X\subseteq
S_0\subsetneq X\cup Y\) and \((X\cup Y)\setminus S_0=N\).  Rigidity of \(Y\)
over \(X\) gives
\[
  |N|-e_G(N;S_0)\alpha<0 .
\]
Every vertex of \(X\cup Y\) lies in \(S\cup N\): vertices of \(X\) and of
\(Y\cap B\) lie in \(B\subseteq S\), and a vertex of \(Y\setminus B=Y'\) lies
in \(S\) or in \(Y'\setminus S=N\).  So an edge with an end in \(N\) and both
ends in \(S_0\cup N=X\cup Y\) has both ends in \(S\cup N\), and is counted by
\(e_G(N;S)\).  Hence \(e_G(N;S)\geq e_G(N;S_0)\) and
\[
  |N|-e_G(N;S)\alpha\leq|N|-e_G(N;S_0)\alpha<0 .
\]
Since \(N=(B\cup Y')\setminus S\), this is the required inequality for \(S\).
\end{proof}
\end{annexitem}

\begin{annexitem}{annex:p03-w03-p03-068-safe-or-rigid}{thm:p03-w03-p03-068-safe-or-rigid}{Safe, or containing a rigid part}
\begin{proof}
For \(U\subseteq Y\) put \(f(U)=|U|-e_G(U;X)\alpha\) for nonempty \(U\) and
\(f(\emptyset)=0\); since \(X\cup Y\) is finite, every \(e_G(U;X)\) is a
finite integer.  The subextensions of \((X,G[X\cup Y])\) are the rooted
graphs \((X,G[X\cup U])\) with \(\emptyset\neq U\subseteq Y\); the nonroots of
such a subextension are the vertices of \(U\), and its counted edges are the
edges of \(G\) with an end in \(U\) and both ends in \(X\cup U\), so its
parameter \(e\) is \(e_G(U;X)\).  Hence \((X,G[X\cup Y])\) is safe exactly
when \(f(U)>0\) for every nonempty \(U\subseteq Y\), and, as \(\alpha\) is
irrational and \(|U|\geq1\), \(f(U)\neq0\) for every nonempty \(U\).

\emph{Not both.}  If \(U\neq\emptyset\) is rigid over \(X\), the case \(S=X\)
of the definition gives \(f(U)<0\), so the subextension with nonroots \(U\) is
dense and \((X,G[X\cup Y])\) is not safe.

\emph{At least one.}  Suppose \((X,G[X\cup Y])\) is not safe, so \(f(U_0)<0\)
for some nonempty \(U_0\subseteq Y\).  Let \(\mu\) be the least value of \(f\)
on the finitely many subsets of \(Y\); then \(\mu\leq f(U_0)<0=f(\emptyset)\).
Choose \(U\) with \(f(U)=\mu\) and \(|U|\) least; \(U\neq\emptyset\).  Let
\(S\) satisfy \(X\subseteq S\subsetneq X\cup U\) and put \(U_1=S\setminus X\),
a proper subset of \(U\), so \((X\cup U)\setminus S=U\setminus U_1\).  The
edges counted by \(e_G(U\setminus U_1;S)\) are the edges with both ends in
\(X\cup U\) and an end in \(U\setminus U_1\).  The edges counted by
\(e_G(U;X)\) are those with both ends in \(X\cup U\) and an end in \(U\);
removing from them those with an end in \(U\setminus U_1\) leaves the edges
with both ends in \(X\cup U_1\) and an end in \(U_1\), which are the edges
counted by \(e_G(U_1;X)\) (none if \(U_1=\emptyset\)).  So
\(e_G(U\setminus U_1;S)=e_G(U;X)-e_G(U_1;X)\), with
\(e_G(\emptyset;X)\) read as \(0\), and
\[
  |U\setminus U_1|-e_G(U\setminus U_1;S)\alpha=f(U)-f(U_1).
\]
Now \(f(U_1)\geq\mu\), and \(f(U_1)=\mu\) is impossible because \(|U_1|<|U|\)
and \(U\) has least size among minimizers.  So \(f(U_1)>f(U)\), and the
displayed number is negative.  As \(S\) was arbitrary, \(U\) is rigid over
\(X\).

\emph{The consequence.}  By definition, \(Z\) is \(t\)-closed exactly when no
\(Y\) with \(1\leq|Y|\leq t\) and \(Y\cap Z=\emptyset\) is rigid over \(Z\).
If every such \((Z,G[Z\cup Y])\) is safe, then by the first part no such
\(Y\) is rigid over \(Z\) (a rigid \(Y\) is its own rigid part), so \(Z\) is
\(t\)-closed.  If some such \((Z,G[Z\cup Y])\) is not safe, the first part
gives a nonempty \(U\subseteq Y\) rigid over \(Z\); then
\(1\leq|U|\leq t\) and \(U\cap Z=\emptyset\), so \(Z\) is not \(t\)-closed.
\end{proof}
\end{annexitem}

\begin{annexitem}{annex:p03-w03-p03-068-well-defined}{thm:p03-w03-p03-068-well-defined}{The two definitions are well defined and agree}
\begin{proof}
\emph{Item 1.}  The set of all vertices is \(t\)-closed, since no nonempty
\(Y\) is disjoint from it.  Let \(\mathcal Z\) be a nonempty family of
\(t\)-closed sets with intersection \(Z^*\), and suppose \(Y\), with
\(1\leq|Y|\leq t\) and \(Y\cap Z^*=\emptyset\), is rigid over \(Z^*\).  If
\(Y\subseteq Z\) for every \(Z\in\mathcal Z\), then \(Y\subseteq Z^*\), which
is impossible for a nonempty \(Y\) disjoint from \(Z^*\).  So there is
\(Z\in\mathcal Z\) with \(Y\setminus Z\neq\emptyset\).  By
\cref{thm:p03-w03-p03-068-move-base} with \(X=Z^*\) and \(B=Z\), the set
\(Y\setminus Z\) is rigid over \(Z\); it is disjoint from \(Z\) and has between
\(1\) and \(t\) elements.  This contradicts the \(t\)-closedness of \(Z\).  So
\(Z^*\) is \(t\)-closed.  The family of \(t\)-closed sets containing \(X\) is
nonempty, its intersection is \(t\)-closed and contains \(X\), and it is
contained in every member: it is the least \(t\)-closed set containing \(X\).

\emph{Item 2.}  Write \(C=\mathrm{cl}_t(X)\) and let
\(X_0\subsetneq X_1\subsetneq\cdots\) be a rigid \(t\)-chain from \(X\).
\(X_0=X\subseteq C\).  Suppose \(X_i\subseteq C\) and
\(X_{i+1}=X_i\cup Y_i\).  If \(Y_i\not\subseteq C\), then by
\cref{thm:p03-w03-p03-068-move-base} with \(X=X_i\) and \(B=C\), the set
\(Y_i\setminus C\) is rigid over \(C\), disjoint from \(C\), with between \(1\)
and \(t\) elements, contradicting item~1.  So \(X_{i+1}\subseteq C\).  By
induction every set of the chain lies in \(C\).

\emph{Item 3.}  The last set \(X_l\) of a terminal chain is \(t\)-closed and
contains \(X\), so \(C\subseteq X_l\) by item~1; and \(X_l\subseteq C\) by
item~2.  The set \(C\) is defined without reference to any chain, so the
outcome does not depend on the choices made.

\emph{Item 4.}  Suppose \(C\setminus X\) has \(k\) elements.  Each step of a
chain adds at least one vertex, which lies in \(C\setminus X\) by item~2, so a
chain has at most \(k\) steps.  A chain ending at \(X_i\) is terminal if
\(X_i\) is \(t\)-closed; otherwise some \(Y\) with \(1\leq|Y|\leq t\),
disjoint from \(X_i\), is rigid over \(X_i\), and the chain extends by
\(X_{i+1}=X_i\cup Y\).  Since lengths are bounded by \(k\), repeating this
reaches a terminal chain, whose last set is \(C\) by item~3.  If \(G\) is
finite, \(C\setminus X\) is finite.

\emph{Item 5.}  \(X\subseteq C\) by definition.  If \(X\subseteq X'\), then
\(\mathrm{cl}_t(X')\) is \(t\)-closed and contains \(X\), so it contains
\(\mathrm{cl}_t(X)\).  If \(s\leq t\), a \(t\)-closed set is \(s\)-closed,
because every \(Y\) with \(1\leq|Y|\leq s\) also has \(|Y|\leq t\); so
\(\mathrm{cl}_t(X)\) is an \(s\)-closed set containing \(X\) and contains
\(\mathrm{cl}_s(X)\).  Finally \(C\) is a \(t\)-closed set containing \(C\),
so \(\mathrm{cl}_t(C)\subseteq C\subseteq\mathrm{cl}_t(C)\).

\emph{Item 6.}  Let \(C\subseteq W\) and let \(D\) be the \(t\)-closure of
\(X\) computed in \(G[W]\).  By the locality part of
\cref{thm:p03-w03-p03-068-rigid-local}, for \(Y\subseteq W\setminus C\) the
set \(Y\) is rigid over \(C\) in \(G[W]\) if and only if it is rigid over \(C\)
in \(G\).  Every candidate \(Y\) in \(G[W]\) is a subset of \(W\), so \(C\) is
\(t\)-closed in \(G[W]\), and \(D\subseteq C\).  Conversely, suppose \(Y\) with
\(1\leq|Y|\leq t\) and \(Y\cap D=\emptyset\) is rigid over \(D\) in \(G\).  If
\(Y\not\subseteq C\), then \(Y\setminus C\) is rigid over \(C\) in \(G\) by
\cref{thm:p03-w03-p03-068-move-base}, contradicting the \(t\)-closedness of
\(C\).  So \(Y\subseteq C\subseteq W\), and by locality \(Y\) is rigid over
\(D\) in \(G[W]\), contradicting the \(t\)-closedness of \(D\) in \(G[W]\).
Hence \(D\) is \(t\)-closed in \(G\); it contains \(X\), so \(C\subseteq D\).
Thus \(D=C\).  Taking \(W=C\) gives the last sentence.
\end{proof}
\end{annexitem}

\begin{annexitem}{annex:p03-w03-p03-068-type-formula}{thm:p03-w03-p03-068-type-formula}{The closure type is first-order definable}
\begin{proof}
\emph{The formula.}  Let \(U=\rho(\{1,\ldots,r\})\) and enumerate
\(V(H)\setminus U\) as \(h_1,\ldots,h_m\).  Assign to each variable a vertex
of \(H\): \(\tau(x_i)=\rho(i)\) and \(\tau(y_j)=h_j\).  Let
\(\Delta_{H,\rho}(\bar x,\bar y)\) be the conjunction of: \(x_i=x_{i'}\) when
\(\rho(i)=\rho(i')\) and \(x_i\neq x_{i'}\) otherwise; \(y_j\neq y_{j'}\) for
\(j\neq j'\); \(x_i\neq y_j\) for all \(i,j\); and, for any two variables
\(u,u'\) with \(\tau(u)\neq\tau(u')\), the formula \(u\sim u'\) if
\(\tau(u)\tau(u')\in E(H)\) and \(\neg\,u\sim u'\) otherwise.

An \emph{attachment of size \(w\)} to \(V(H)\), for \(1\leq w\leq t\), is a set
\(P\) of unordered pairs from \(V(H)\sqcup\{1,\ldots,w\}\), each pair having at
least one element in \(\{1,\ldots,w\}\).  Let \(H\cup P\) be the graph on
\(V(H)\sqcup\{1,\ldots,w\}\) with edge set \(E(H)\cup P\).  The attachment is
\emph{rigid} if \(\{1,\ldots,w\}\) is rigid over \(V(H)\) in \(H\cup P\).  Every
edge counted in this rigidity condition has an end in \(\{1,\ldots,w\}\), so it
lies in \(P\); rigidity depends only on \(P\) and \(\alpha\).  There are
finitely many attachments with \(w\leq t\).  For an attachment \(P\) of size
\(w\) let \(\Theta_P(\bar x,\bar y)\) be
\[
  \exists z_1\cdots\exists z_w\Bigl[\bigwedge_{k<k'}z_k\neq z_{k'}\wedge
  \bigwedge_{k,u}z_k\neq u\wedge\bigwedge_{\{a,b\}}\bigl(\pm\,
  \sigma(a)\sim\sigma(b)\bigr)\Bigr],
\]
where \(u\) ranges over the variables \(\bar x,\bar y\); \(\{a,b\}\) ranges
over the pairs from \(V(H)\sqcup\{1,\ldots,w\}\) with at least one element in
\(\{1,\ldots,w\}\); \(\sigma(k)=z_k\) and \(\sigma(h)\) is a variable \(u\) with
\(\tau(u)=h\) (for a root vertex take the least \(i\) with \(\rho(i)=h\)); and
\(\pm\) is no sign if \(\{a,b\}\in P\) and \(\neg\) otherwise.  Let
\(\Theta_{t,H}\) be the disjunction of \(\Theta_P\) over the rigid
attachments \(P\), and \(\varphi_{t,H,\rho}=\exists\bar
y\,[\Delta_{H,\rho}\wedge\neg\Theta_{t,H}]\).

\emph{Claim A.}  If \(\Delta_{H,\rho}\) holds in \(G\) at \(\bar a,\bar b\), let
\(W=\{a_1,\ldots,a_r,b_1,\ldots,b_m\}\) and define \(\iota\colon V(H)\to W\) by
\(\iota(\rho(i))=a_i\) and \(\iota(h_j)=b_j\).  It is well defined because
\(\rho(i)=\rho(i')\) forces \(a_i=a_{i'}\); it is injective by the inequality
clauses and onto by construction; and it preserves adjacency and nonadjacency
by the remaining clauses.  So \(\iota\) is an isomorphism of \(H\) onto
\(G[W]\) with \(\iota(\rho(i))=a_i\).

\emph{Claim B.}  In the situation of Claim A, \(\Theta_{t,H}\) holds at
\(\bar a,\bar b\) if and only if \(W\) is not \(t\)-closed in \(G\).  If
\(\Theta_P\) holds with witnesses \(c_1,\ldots,c_w\), then
\(Y=\{c_1,\ldots,c_w\}\) is disjoint from \(W\), has \(w\) elements, and
\(\iota\) extended by \(k\mapsto c_k\) is an isomorphism of \(H\cup P\) onto
\(G[W\cup Y]\), by Claim A and the exact adjacency clauses of \(\Theta_P\).
Rigidity is preserved by isomorphism, so \(Y\) is rigid over \(W\) in
\(G[W\cup Y]\), hence in \(G\) by the locality part of
\cref{thm:p03-w03-p03-068-rigid-local}; so \(W\) is not \(t\)-closed.
Conversely, if \(Y=\{c_1,\ldots,c_w\}\), \(1\leq w\leq t\), is disjoint from
\(W\) and rigid over \(W\), let \(P\) be the set of pairs \(\{a,b\}\) with at
least one element in \(\{1,\ldots,w\}\) whose images under
\(\iota\cup(k\mapsto c_k)\) are adjacent in \(G\).  Then that map is an
isomorphism of \(H\cup P\) onto \(G[W\cup Y]\), \(P\) is rigid by locality and
invariance under isomorphism, and \(c_1,\ldots,c_w\) witness \(\Theta_P\).

\emph{Correctness.}  Let \(X=\{x_1,\ldots,x_r\}\) and \(C=\mathrm{cl}_t(X)\).
Suppose first \(\mathrm{cl}_t(\bar x)\cong H\) by an isomorphism
\(\pi\colon G[C]\to H\) with \(\pi(x_i)=\rho(i)\).  Put
\(b_j=\pi^{-1}(h_j)\).  Then \(\pi(X)=U\), so the \(b_j\) are distinct from
the \(x_i\); \(\Delta_{H,\rho}\) holds at \(\bar x,\bar b\) because \(\pi\) is an
isomorphism; \(W=C\) is \(t\)-closed, so \(\neg\Theta_{t,H}\) holds by Claim B.
Hence \(G\entails\varphi_{t,H,\rho}[\bar x]\).

Suppose next that \(\bar b\) witnesses \(\Delta_{H,\rho}\wedge\neg\Theta_{t,H}\)
at \(\bar x\).  By Claim A there is an isomorphism \(\iota\colon H\to G[W]\)
with \(\iota(U)=X\), and by Claim B the set \(W\) is \(t\)-closed in \(G\), so
\(C\subseteq W\).  By item~6 of \cref{thm:p03-w03-p03-068-well-defined}, the
\(t\)-closure of \(X\) computed in \(G[W]\) is \(C\).  The definitions of
rigidity and closure are invariant under isomorphism, so that closure is
\(\iota(\mathrm{cl}^H_t(U))=\iota(V(H))=W\), using that \((H,\rho)\) is a
\(t\)-closure type.  So \(C=W\), and \(\iota^{-1}\colon G[C]\to H\) is an
isomorphism with \(x_i\mapsto\rho(i)\).

\emph{The case \(t=0\).}  Every set is \(0\)-closed, so the \(0\)-closure of
\(U\) in \(H\) is \(U\), whence \(V(H)=U\) and \(m=0\); there are no
attachments, \(\Theta_{0,H}\) is the empty disjunction (false), and
\(\varphi_{0,H,\rho}\) is equivalent to the quantifier-free
\(\Delta_{H,\rho}(\bar x)\).  It says that \(x_i\mapsto\rho(i)\) is a well-defined
isomorphism of \(G[X]\) onto \(H\), matching
\cref{def:p03-w03-p03-068-closure-iso}.

\emph{Why the type condition is needed.}  Let \(t\geq1\), \(r=1\), and let
\(H\) have two vertices \(\rho(1),h_1\) and no edge.  Then \(\mathrm{cl}^H_t(\{\rho(1)\})=\{\rho(1)\}\)
because one vertex with no counted edge gives \(1-0\cdot\alpha>0\); so \((H,\rho)\)
is not a \(t\)-closure type.  In the graph \(G\) with two vertices \(x,b\) and
no edge, \(\Delta_{H,\rho}\wedge\neg\Theta_{t,H}\) holds at \(x,b\), because
\(\{x,b\}\) is the whole vertex set and so is \(t\)-closed; but
\(\mathrm{cl}_t(x)=\{x\}\) is not isomorphic to \(H\).

The formula depends on \(\alpha\) only through which attachments are rigid.
\end{proof}
\end{annexitem}

\begin{annexitem}{annex:p03-w03-p03-068-no-dense}{thm:p03-w03-p03-068-no-dense}{No dense set of bounded size}
\begin{proof}
Let \(p=n^{-\alpha}\leq1\).  For \(1\leq s\leq M\) let \(e_s\) be the least
integer with \(s-e_s\alpha<0\); it exists because \(\alpha>0\).  A set \(S\)
with \(|S|=s\) is dense exactly when \(e(S)\geq e_s\), and then \(G[S]\)
contains a set \(F\) of exactly \(e_s\) edges.  If \(e_s>\binom s2\) there is
no such \(S\).  For a fixed \(s\)-set \(S\) and a fixed set \(F\) of \(e_s\)
pairs inside \(S\), the probability that every pair in \(F\) is an edge is
\(p^{e_s}\), by independence of the edges.  The union bound over the
\(\binom ns\leq n^s\) sets \(S\) and the \(\binom{\binom s2}{e_s}\) sets \(F\)
gives
\begin{align*}
  &\Pr[\text{some dense set of at most }M\text{ vertices}]\\
  &\qquad\leq\sum_{s=1}^{M}\binom{\binom s2}{e_s}n^{s}n^{-\alpha e_s}
  =\sum_{s=1}^{M}\binom{\binom s2}{e_s}n^{s-e_s\alpha},
\end{align*}
with the convention that a binomial coefficient \(\binom ab\) with \(b>a\) is
\(0\).  Each of the \(M\) exponents \(s-e_s\alpha\) is negative and the
coefficients do not depend on \(n\), so the sum tends to \(0\).
\end{proof}
\end{annexitem}

\begin{annexitem}{annex:p03-w03-p03-068-nonexistence}{thm:p03-w03-p03-068-nonexistence}{Nonexistence Lemma}
\begin{proof}
\(\mathrm{cl}_t(\emptyset)=\emptyset\) holds exactly when \(\emptyset\) is
\(t\)-closed: if \(\emptyset\) is \(t\)-closed it is the least \(t\)-closed set
containing \(\emptyset\), and conversely \(\mathrm{cl}_t(\emptyset)\) is
\(t\)-closed by item~1 of \cref{thm:p03-w03-p03-068-well-defined}.  And
\(\emptyset\) is \(t\)-closed exactly when no set \(Y\) with
\(1\leq|Y|\leq t\) is rigid over \(\emptyset\).  This proves the first
equivalence.

If \(Y\) is rigid over \(\emptyset\), the case \(S=\emptyset\) of the
definition gives \(|Y|-e_G(Y;\emptyset)\alpha<0\), and
\(e_G(Y;\emptyset)=e(Y)\); so \(Y\) is a dense set of at most \(t\) vertices.
By \cref{thm:p03-w03-p03-068-no-dense} with \(M=t\), almost surely there is no
such set, so almost surely \(\mathrm{cl}_t(\emptyset)=\emptyset\).

If \(H\) has \(v\) vertices, \(e\) edges and \(v-e\alpha<0\), the vertex set
\(S\) of a copy of \(H\) in \(G\) has \(|S|=v\) and \(e(S)\geq e\), so
\(|S|-e(S)\alpha\leq v-e\alpha<0\) and \(S\) is dense.  By
\cref{thm:p03-w03-p03-068-no-dense} with \(M=v\), almost surely there is no
copy of \(H\).
\end{proof}
\end{annexitem}

\begin{annexitem}{annex:p03-w03-p03-068-finite-closure}{thm:p03-w03-p03-068-finite-closure}{Finite Closure Lemma}
\begin{proof}
\emph{The constant \(\epsilon\).}  Fix \(v\) with \(1\leq v\leq t\).  Since
\(\alpha\) is irrational, \(v/\alpha\) is irrational, so the integers \(e\)
with \(v-e\alpha\leq0\) are those with \(e\geq e_v=\lceil v/\alpha\rceil\), and
\(e_v\alpha>v\).  The quantity \((e\alpha-v)/v\) increases with \(e\), so its
minimum over these \(e\) is \((e_v\alpha-v)/v>0\).  Hence \(\epsilon\) is the
minimum of the \(t\) positive numbers \((e_v\alpha-v)/v\), \(1\leq v\leq t\);
it exists and is positive.

\emph{\(K>r/\epsilon\).}  Let \(\epsilon=(e_v\alpha-v)/v\).  If \(r/\epsilon\)
were a rational \(c\), then \(c>0\) and \(e_v\alpha=v+rv/c\) would be rational,
and since \(e_v\geq1\) so would \(\alpha\).  So \(r/\epsilon\) is irrational,
\(K=\lceil r/\epsilon\rceil>r/\epsilon\), \(K\geq1\), and \(r-K\epsilon<0\).

\emph{A deterministic statement.}  Put \(M=r+K+t-1\).  We show: if a finite
graph \(G\) has no dense set of at most \(M\) vertices, then
\(|\mathrm{cl}_t(x_1,\ldots,x_r)|<r+K\) for all vertices \(x_1,\ldots,x_r\).
Together with \cref{thm:p03-w03-p03-068-no-dense}, which says that almost
surely \(G(n,n^{-\alpha})\) has no such dense set, this proves the lemma.

Suppose instead that \(X=\{x_1,\ldots,x_r\}\), with \(r_0=|X|\leq r\), has
\(C=\mathrm{cl}_t(X)\) of size at least \(r+K\).  By item~4 of
\cref{thm:p03-w03-p03-068-well-defined} there is a terminal chain
\(X=X_0\subsetneq\cdots\subsetneq X_l=C\).  Since \(|X_0|=r_0<r+K\), there is a
least \(j\geq1\) with \(|X_j|\geq r+K\); then \(|X_{j-1}|\leq r+K-1\) and
\(|X_j|\leq|X_{j-1}|+t\leq M\).  Put \(W=X_j\) and, for \(1\leq i\leq j\),
\(N_i=X_i\setminus X_{i-1}\), \(v_i=|N_i|\in[1,t]\) and
\(e_i=e_G(N_i;X_{i-1})\).

Each \(N_i\) is rigid over \(X_{i-1}\); the case \(S=X_{i-1}\) gives
\(v_i-e_i\alpha<0\), so \(e_i\) is one of the integers in the definition of
\(\epsilon\) for \(v=v_i\), and \((e_i\alpha-v_i)/v_i\geq\epsilon\), that is,
\[
  v_i-e_i\alpha\leq-\epsilon v_i .
\]
An edge counted by \(e_i\) has both ends in \(X_i\).  An edge counted by
\(e_k\) with \(k>i\) has an end in \(N_k\), which is disjoint from \(X_i\).
So no edge is counted twice, and all counted edges lie in \(G[W]\); hence
\(e(W)\geq\sum_{i\leq j}e_i\).  Also \(|W|=r_0+\sum_{i\leq j}v_i\), and
\(\sum_{i\leq j}v_i=|W|-r_0\geq r+K-r_0\geq K\).  Therefore
\[
  |W|-e(W)\alpha\leq r_0+\sum_{i=1}^{j}(v_i-e_i\alpha)
  \leq r_0-\epsilon\sum_{i=1}^{j}v_i\leq r-\epsilon K<0 .
\]
So \(W\) is a dense set of at most \(M\) vertices, a contradiction.
\end{proof}
\end{annexitem}

\begin{annexitem}{annex:p03-w03-p03-068-finitely-many-types}{thm:p03-w03-p03-068-finitely-many-types}{Finitely many closure types occur}
\begin{proof}
Let \(K\) be as in \cref{thm:p03-w03-p03-068-finite-closure}.  Up to
isomorphisms respecting \(\rho\) there are finitely many pairs \((H,\rho)\)
with \(|V(H)|<r+K\) and \(\rho\colon\{1,\ldots,r\}\to V(H)\); let \(\mathcal C\)
contain one representative of each class that is a \(t\)-closure type.
Almost surely every \(r\)-tuple \(\bar x\) has
\(C=\mathrm{cl}_t(\bar x)\) with \(|C|<r+K\).  By item~6 of
\cref{thm:p03-w03-p03-068-well-defined}, the \(t\)-closure of
\(\{x_1,\ldots,x_r\}\) computed in \(G[C]\) is \(C\), so \(G[C]\) with
\(\rho(i)=x_i\) is a \(t\)-closure type with fewer than \(r+K\) vertices, and it
is isomorphic, by a map sending each \(x_i\) to \(\rho(i)\), to a member of
\(\mathcal C\).  By \cref{thm:p03-w03-p03-068-type-formula}, \(\bar x\)
satisfies the corresponding \(\varphi_{t,H,\rho}\).
\end{proof}
\end{annexitem}

\begin{annexitem}{annex:p03-w03-p03-068-t-fixed}{thm:p03-w03-p03-068-t-fixed}{The bound \(t\) must be fixed}
\begin{proof}
\emph{A dense set contains a rigid one.}  Let \(W\) be a dense set in a finite
graph \(G\).  The rooted graph \((\emptyset,G[W])\) has the subextension
\((\emptyset,G[W])\) itself, which is dense, so it is not safe.  By
\cref{thm:p03-w03-p03-068-safe-or-rigid} with \(X=\emptyset\) and \(Y=W\),
some nonempty \(U\subseteq W\) is rigid over \(\emptyset\).

\emph{The whole vertex set is almost surely dense.}  The number \(E\) of edges
of \(G(n,p)\), \(p=n^{-\alpha}\), is a sum of \(\binom n2\) independent
indicators of probability \(p\), with mean \(\mu_n=\binom n2p\) and variance
\(\binom n2p(1-p)\leq\mu_n\).  Applying \(\Pr[Z\geq c]\leq\mathbb E[Z]/c\) to
\(Z=(E-\mu_n)^2\) and \(c=\mu_n^2/4\) gives
\(\Pr[|E-\mu_n|\geq\mu_n/2]\leq4/\mu_n\).  For \(n\geq2\),
\(\mu_n=n(n-1)n^{-\alpha}/2\geq n^{2-\alpha}/4\), which tends to infinity, so
almost surely \(E>\mu_n/2\geq n^{2-\alpha}/8\).  Then
\(n-E\alpha<n-\alpha n^{2-\alpha}/8=n(1-\alpha n^{1-\alpha}/8)\), which is
negative once \(n^{1-\alpha}>8/\alpha\), since \(1-\alpha>0\).  So almost
surely \(V(G)\) is dense.

\emph{Conclusion.}  On that event, the first part gives a nonempty
\(U\subseteq V(G)\) rigid over \(\emptyset\), with \(|U|\leq n\).  So
\(\emptyset\) is not \(n\)-closed, and
\(\mathrm{cl}_n(\emptyset)\neq\emptyset\).
\end{proof}
\end{annexitem}

\begin{annexitem}{annex:p03-w03-p03-068-axiomatized}{thm:p03-w03-p03-068-axiomatized}{Almost sure theory, axiomatized form}
\begin{proof}
Let \(\Sigma_\alpha\) contain the graph axioms and the Nonexistence and
Generic Extension axioms of \cref{def:p03-w03-p03-068-axioms}.  A copy of
a graph here is an injective edge-preserving map; it need not reflect
nonedges.  These are exactly the schemas compared, axiom by axiom, in
\cref{thm:p03-w03-p03-067-axioms-generate}(1).  Item~(2) of that lemma
gives almost sure validity of each member of \(\Sigma_\alpha\), and
item~(3) gives consistency, completeness and, for every graph sentence
\(A\),
\[
 A\in T_\alpha\quad\Longleftrightarrow\quad\Sigma_\alpha\entails A,
\]
where \(T_\alpha\) is the set of almost sure sentences defined in
\cref{def:p03-w03-p03-068-fa}.  This is the claimed axiomatization.
The cited lemma obtains its probabilistic input from
\cref{thm:p03-w03-p03-070-as-holds} and completeness from
\cref{thm:p03-w03-p03-070-complete}; it does not assume the present
axiomatization or the continuity theorem below.
\end{proof}
\end{annexitem}

\begin{annexitem}{annex:p03-w03-p03-068-continuity}{thm:p03-w03-p03-068-continuity}{Continuity of \(f_A\)}
\begin{proof}
\emph{Step 1: signs persist.}  For integers \(v\) and \(e\geq1\),
\(v-e\alpha_0\neq0\) because \(\alpha_0\) is irrational; put
\(\delta(v,e)=|v-e\alpha_0|/e>0\).  If \(|\alpha-\alpha_0|<\delta(v,e)\), then
\(|(v-e\alpha)-(v-e\alpha_0)|=e|\alpha-\alpha_0|<|v-e\alpha_0|\), so
\(v-e\alpha\) has the sign of \(v-e\alpha_0\).  For \(e=0\) the number \(v\)
does not depend on \(\alpha\).

\emph{Step 2: each axiom persists.}  Let \(\sigma\) be an axiom of
\(T_{\alpha_0}\) in the sense of \cref{def:p03-w03-p03-068-axioms},
including the graph axioms.  The graph axioms do not depend on the
parameter and persist with any positive bound.  If
\(\sigma\) is the Nonexistence axiom for \(H\), with \(v\) vertices and \(e\)
edges, then \(v-e\alpha_0<0\) forces \(e\geq1\); for
\(|\alpha-\alpha_0|<\delta(v,e)\) also \(v-e\alpha<0\), and the sentence
``there is no copy of \(H\)'' does not mention \(\alpha\), so \(\sigma\) is an
axiom of \(T_\alpha\).  If \(\sigma\) is the Generic Extension axiom for a safe
rooted graph \((R,H)\) and \(t\), the sentence is determined by \((R,H)\), \(t\)
and the set of rigid attachments of at most \(t\) new vertices to \(r+v\)
labelled vertices that have an edge to a nonroot (as in the proof of
\cref{thm:p03-w03-p03-068-type-formula}).  Whether \((R,H)\) is safe depends on
the signs of \(v'-e'\alpha\) for its finitely many subextensions; whether each
of the finitely many attachments is rigid depends on the signs of
\(v'-e'\alpha\) for its finitely many nailextensions.  Let \(F_\sigma\) be the
finite set of pairs \((v',e')\) involved and
\(\delta_\sigma=\min\{\delta(v',e'):(v',e')\in F_\sigma,\ e'\geq1\}\) (or
\(1\) if there is none).  For \(|\alpha-\alpha_0|<\delta_\sigma\), Step~1 shows
that \((R,H)\) is safe at \(\alpha\) and that the set of rigid attachments is
the same, so \(\sigma\) is literally the Generic Extension axiom for \((R,H)\)
and \(t\) at \(\alpha\).

\emph{Step 3.}  Since \(\Pr[\neg A]=1-\Pr[A]\), \(f_{\neg A}=1-f_A\); so we may
assume \(f_{A}(\alpha_0)=1\), replacing \(A\) by \(\neg A\) otherwise.  Then
\(A\in T_{\alpha_0}\), and by \cref{thm:p03-w03-p03-068-axiomatized} \(A\) is a
logical consequence of the axioms of \(T_{\alpha_0}\).  By the compactness
theorem (\cref{thm:fol-compactness}) it is a consequence of a finite set
\(\Sigma_0\) of them.  Let \(\delta>0\) be the minimum of the \(\delta_\sigma\),
\(\sigma\in\Sigma_0\), and of \(\min(\alpha_0,1-\alpha_0)\).  For irrational
\(\alpha\) with \(|\alpha-\alpha_0|<\delta\), every member of \(\Sigma_0\) is
an axiom of \(T_\alpha\) by Step~2, hence holds almost surely in
\(G(n,n^{-\alpha})\) by \cref{thm:p03-w03-p03-068-axiomatized}, hence lies in
\(T_\alpha\).  Since \(T_\alpha\) is closed under logical consequence
(\cref{thm:p03-w03-p03-067-as-theory}), \(A\in T_\alpha\) and
\(f_A(\alpha)=1=f_A(\alpha_0)\).
\end{proof}
\end{annexitem}

\begin{annexitem}{annex:p03-w03-p03-068-distinct}{thm:p03-w03-p03-068-distinct}{The theories \(T_\alpha\) are pairwise distinct}
\begin{proof}
\emph{Writing \(q\).}  Since \(1/q>1\) is rational, write
\(1/q=(d+1)+j/a\) with integers \(d=\lfloor1/q\rfloor-1\geq0\), \(a\geq3\) and
\(0\leq j<a\): if the fractional part of \(1/q\) is \(j_0/a_0\) in lowest terms,
take \(a=3a_0\) and \(j=3j_0\).

\emph{Even spreading.}  For every integer \(i\) put
\(h_i=\lfloor(i+1)j/a\rfloor-\lfloor ij/a\rfloor\).  Since \(0\leq j/a<1\),
\(h_i\in\{0,1\}\); since \(\lfloor(i+a)j/a\rfloor=\lfloor ij/a\rfloor+j\),
\(h_{i+a}=h_i\) and \(\sum_{i=0}^{a-1}h_i=j\).  For every integer \(s\) and
\(L\geq1\) the sum telescopes:
\[
  \sum_{i=s}^{s+L-1}h_i=\Bigl\lfloor\frac{(s+L)j}a\Bigr\rfloor
  -\Bigl\lfloor\frac{sj}a\Bigr\rfloor<\frac{Lj}a+1,
\]
because \(\lfloor x+y\rfloor\leq x+y\) and \(\lfloor x\rfloor>x-1\).

\emph{The rooted graph.}  \((R_q,H_q)\) has roots \(c_0,\ldots,c_d\) and
nonroots \(y_0,\ldots,y_{a-1}\).  Its edges are the cycle edges
\(y_iy_{i+1}\) (indices modulo \(a\); a cycle of length \(a\geq3\)), the edges
\(y_ic_k\) for all \(i\) and \(k<d\), and the edge \(y_ic_d\) exactly when
\(h_i=1\).  There are no edges between roots.  For a nonempty set
\(N\) of nonroots, the subextension with nonroot set \(N\) has
\(|N|\) nonroots and
\[
  e(N)=\sum_{y_i\in N}(d+h_i)+c(N)
\]
counted edges, where \(c(N)\) is the number of cycle edges with both ends in
\(N\).  For \(N\) the set of all nonroots, \(e(N)=da+j+a=a/q\).

\emph{Claim: \(e(N)<|N|/q\) for every nonempty proper subset \(N\).}  Split
\(N\) into its maximal runs of cyclically consecutive indices, \(k\geq1\) arcs
of lengths \(L_1,\ldots,L_k\).  Each arc has length less than \(a\), so it
contains exactly \(L_m-1\) cycle edges, and no cycle edge joins two different
maximal runs; hence \(c(N)=|N|-k\).  By even spreading, the arc of length
\(L_m\) contains fewer than \(L_mj/a+1\) indices with \(h_i=1\).  So
\[
  e(N)<d|N|+\Bigl(\frac{|N|j}a+k\Bigr)+(|N|-k)=|N|\Bigl(d+1+\frac ja\Bigr)
  =\frac{|N|}q .
\]

\emph{The two regimes.}  Let \(\alpha<q\).  For every nonempty \(N\),
\(e(N)\leq|N|/q\), so \(|N|-e(N)\alpha\geq|N|(1-\alpha/q)>0\): every
subextension is sparse, \((R_q,H_q)\) is safe, and
\(\operatorname{Ext}(R_q,H_q)\) holds almost surely by
\cref{thm:p03-w03-p03-065-safe-criterion}.  Let \(\alpha>q\).  For \(N\) the
set of all nonroots, \(a-(a/q)\alpha=a(1-\alpha/q)<0\), so the subextension
\((R_q,H_q)\) itself is not sparse, \((R_q,H_q)\) is not safe, and
\(\operatorname{Ext}(R_q,H_q)\) holds almost never by the same theorem.

\emph{Distinctness.}  Given irrational \(\alpha<\beta\) in \((0,1)\), choose a
rational \(q\) with \(\alpha<q<\beta\).  Then \(\operatorname{Ext}(R_q,H_q)\in
T_\alpha\) and \(\neg\operatorname{Ext}(R_q,H_q)\in T_\beta\); since
\(T_\beta\) is consistent (\cref{thm:p03-w03-p03-067-as-theory}),
\(\operatorname{Ext}(R_q,H_q)\notin T_\beta\).
\end{proof}
\end{annexitem}

\begin{annexitem}{annex:p03-w03-p03-068-clique-spectra}{thm:p03-w03-p03-068-clique-spectra}{Spectra from cliques}
\begin{proof}
\emph{Boolean combinations.}  For sentences \(B,C\) and irrational \(\alpha\),
\(f_{\neg B}=1-f_B\), and \(f_{B\wedge C}=f_Bf_C\): if both limits are \(1\),
\(\Pr[\neg(B\wedge C)]\leq\Pr[\neg B]+\Pr[\neg C]\to0\); if one is \(0\),
\(\Pr[B\wedge C]\) is at most that probability and tends to \(0\).

\emph{One clique.}  Put \(c_k=2/(k-1)\).  For \(p=n^{-\alpha}\),
\(p/n^{-c_k}=n^{c_k-\alpha}\) tends to infinity if \(\alpha<c_k\) and to \(0\)
if \(\alpha>c_k\).  By \cref{thm:p03-w03-p03-065-clique}, \(A_k\) holds almost
surely in the first case and almost never in the second.  So \(f_{A_k}\) is
\(1\) on the irrationals below \(c_k\) and \(0\) on those above.

\emph{Spectra.}  For \(k=3\), \(c_3=1\), so \(f_{A_3}\equiv1\) on \((0,1)\)
and \(\operatorname{Sp}(A_3)=\emptyset\).  For \(k=4\), \(c_4=2/3\): every interval around
\(2/3\) contains irrationals on both sides of \(2/3\) inside \((0,1)\), where
\(f_{A_4}\) takes the values \(1\) and \(0\); so \(2/3\in \operatorname{Sp}(A_4)\).  Any other
\(q\in(0,1)\) has an interval around it avoiding \(2/3\), on which
\(f_{A_4}\) is constant.  So \(\operatorname{Sp}(A_4)=\{2/3\}\), and in the same way
\(\operatorname{Sp}(A_5)=\{1/2\}\).  By the first paragraph
\(f_{A_4\wedge\neg A_5}=f_{A_4}(1-f_{A_5})\), which is \(1\) on the irrationals
of \((1/2,2/3)\) and \(0\) on those of \((0,1/2)\cup(2/3,1)\); the same
argument gives \(\operatorname{Sp}(A_4\wedge\neg A_5)=\{1/2,2/3\}\).
\end{proof}
\end{annexitem}

\begin{annexitem}{annex:p03-w03-p03-068-finite-spectra}{thm:p03-w03-p03-068-finite-spectra}{Every finite set of rationals is a spectrum}
\begin{proof}
\emph{Step 1.}  By the first paragraph of the proof of
\cref{thm:p03-w03-p03-068-clique-spectra}, \(f\) of a Boolean combination is
the same Boolean combination of the values \(f\).  In particular, for
\(B\oplus C=(B\wedge\neg C)\vee(\neg B\wedge C)\), \(f_{B\oplus C}=f_B+f_C\)
modulo \(2\).

\emph{Step 2.}  If \(Q=\emptyset\), let \(A\) be \(\exists x\,x=x\); it holds in
every graph with at least one vertex, so \(f_A\equiv1\) and
\(\operatorname{Sp}(A)=\emptyset\).  Otherwise let \(Q=\{q_1<\cdots<q_k\}\) and let
\(E_i=\operatorname{Ext}(R_{q_i},H_{q_i})\) be the sentences of
\cref{thm:p03-w03-p03-068-distinct}, so that \(f_{E_i}(\alpha)=1\) if
\(\alpha<q_i\) and \(0\) if \(\alpha>q_i\).  Let
\(A=E_1\oplus\cdots\oplus E_k\).  By Step~1, \(f_A(\alpha)\) is the parity of
the number of \(i\) with \(\alpha<q_i\).  Put \(q_0=0\), \(q_{k+1}=1\).  On the
irrationals of \((q_m,q_{m+1})\) that number is \(k-m\), a constant.  A point
\(q\in(0,1)\setminus Q\) lies in one of these open intervals, and a small
interval around it lies inside, so \(q\notin \operatorname{Sp}(A)\).  For \(q=q_m\),
\(1\leq m\leq k\), every interval around \(q_m\) contains irrationals of
\((q_{m-1},q_m)\), where the number is \(k-m+1\), and of \((q_m,q_{m+1})\),
where it is \(k-m\); the parities differ, so \(q_m\in \operatorname{Sp}(A)\).  Hence
\(\operatorname{Sp}(A)=Q\).
\end{proof}
\end{annexitem}
